\documentclass[11pt]{article}

\usepackage[a4paper,margin=2.5cm]{geometry}
\usepackage{times}
\usepackage{booktabs}
\usepackage{amsmath}
\usepackage{hyperref}
\usepackage{parskip}
\usepackage{caption}
\usepackage{listings}
\usepackage{xcolor}

\emergencystretch=3em

\lstdefinestyle{md}{
  basicstyle=\ttfamily\scriptsize,
  breaklines=true,
  columns=fullflexible,
  backgroundcolor=\color{gray!8},
  frame=single,
  framerule=0.3pt,
  xleftmargin=2pt,
  xrightmargin=2pt
}
\lstset{style=md}

\title{From Black Box to White Box:\\A Case Study in Autoresearch}
\author{Christian Block\\\small Research project with Prof.\ Ramesh}
\date{July 2026}

\begin{document}
\maketitle

\begin{abstract}
A language model can execute a research loop on its own: given a markdown file of instructions and a way to run experiments, it repeatedly proposes something to try, runs it, looks at the result, and decides what to do next. Change the instructions and the loop behaves differently, but it is unclear which part of the instructions is responsible for which part of the difference. This research project treats the instruction file as a configuration, in the same way an algorithm's hyperparameters form a configuration, and borrows a standard technique from AutoML to measure the effect of each instructional choice on the loop's behavior. Five instructional choices were varied, an eight-configuration ablation design was run against a real machine learning task with an independent language model (OpenAI \texttt{gpt-4o-mini}) acting as the agent on every run, and a simple statistical model was fit to explain which instructional choice drives which aspect of the resulting behavior. The method cleanly recovers which instruction governs each aspect of the agent's \emph{process} each effect traceable to the specific paragraph of the instruction file responsible. It also returns a pointed negative result: none of the five instructional choices reliably moves the accuracy the agent ultimately reaches, which stays near the dataset's ceiling however the instructions are worded. A held-out prediction, made before the corresponding configuration was ever run, is checked against a real execution: the process quantity is predicted almost exactly, while the outcome quantity is correctly flagged as unpredictable.
\end{abstract}

\section{Introduction}

Suppose one wants a language model to act as a autonomous researcher. One gives it a description of a task together with a tool it can call to actually run an experiment and get a result back. The model reads the instructions, proposes a plan to try, calls the tool, reads the result, and decides what to do next, repeating this until it decides to stop. Systems built this way are sometimes called autoresearch loops, and they are attractive precisely because they need so little supervision: point one at a problem, and it works through the same cycle a human researcher would without a person in the loop.

However, the behavior of such a system is guided by one component: the instruction file. We refer to this file throughout as \texttt{program.md}. Anyone who has run one of these loops more than a few times notices that small changes to the wording of the instructions produce noticeably different research behavior. We notice that the instructions matter, but we do not have a way of evaluating which part of the instructions is responsible for which part of the resulting behavior.

This project treats that gap as a measurement problem and borrows the tools for it from AutoML that already solved a structurally identical problem: When a learning algorithm has several hyperparameters, researchers do not simply observe that changing hyperparameters changes performance and leave it there. They routinely apply techniques such as ablation analysis, fractional factorial experiment design, and variance decomposition over hyperparameters to identify which hyperparameters matter, how much each one matters, and whether they interact \cite{fawcetthoos,hutter2014fanova,plackettburman}. Our observation is that \texttt{program.md} is similar to a 
hyperparameter vector.

To keep the experiment compact enough to actually run and check by hand, the domain was restricted to tuning Random Forest and Gradient Boosting classifiers. The point of choosing a well-understood task is that it lets us focus entirely on the question whether changing the instructions change the research process in a way we can measure, attribute, and predict?

The rest of this report is organized as follows. Section~\ref{sec:background} briefly explains the three ideas borrowed from automated machine learning. Section~\ref{sec:methodology} walks through the five steps of the executed research plan. Section~\ref{sec:setup} states the exact experimental setup. Section~\ref{sec:results} presents and discusses the results. Section~\ref{sec:limitations} discusses limitation, and Section~\ref{sec:conclusion} concludes.

\section{Background}
\label{sec:background}

Three ideas underlie this project which were originally developed for understanding algorithms with many hyperparameters rather than instruction files.

\textbf{Variance decomposition over configuration}. If an algorithm has several hyperparameters and its performance varies as we change them, a natural question is how much of that variation each hyperparameter is individually responsible for. Functional ANOVA \cite{hutter2014fanova} answers this by treating the outcome as a function of the hyperparameter vector and decomposing its variance into a part attributable to each hyperparameter on its own, plus parts attributable to interactions between pairs of hyperparameters. The result is a short list of interpretable numbers instead of a black-box performance curve.

\textbf{Ablation analysis}: given two specific configurations with different performance, which of the individual differences between them is responsible for how much of the gap \cite{fawcetthoos}. This complements variance decomposition by giving a step-by-step account of one specific comparison, rather than a global summary across many configurations.

\textbf{Fractional factorial design}: measuring the effect of $k$ two-level choices requires testing all $2^k$ combinations. Fractional factorial design sidesteps this combinatorial explosion by testing a subset of combinations. This design guarantees that the individual effect of every choice can be estimated without bias from the others. The classical construction for this, due to Plackett and Burman \cite{plackettburman}, produces such a subset in multiples of four runs.
\section{Method}
\label{sec:methodology}

The project follows a five-step plan: turn the instruction file into a configuration (Step 1). Decide which configurations to actually run and run them against a real experiment executor (Step 2). Quantify each run (Step 3). Fit a model relating configuration to those numbers (Step 4). And check that model's predictions against a real, held-out run (Step 5).

\subsection{Step 1: turning the instruction file into a configuration}
\label{sec:configspace}

The first executed step reduces the instruction file to a  structured set of choices. Five instructional choices were identified by reading through example instruction files and noting which sentences plausibly influence behavior. Each choice was restricted to exactly two level which results in a a five-bit configuration.

A configuration is then simply a vector of five zeros and ones, for instance $(1,0,0,1,1)$, one bit per axis in the order metric, breadth, stopping rule, output format, emphasis. Any two generated instruction files differ only in the axes that were actually varied between them.

The metric axis controls how precisely the instructions define success. The breadth axis controls whether the agent must commit to a single model family, Random Forest or Gradient Boosting, for an entire run, or is free to try both. The stopping-rule axis controls the experiment budget. The output-format axis controls whether responses must be terse and numeric, or must include full, reasoned explanations of hypotheses and interpretations. The emphasis axis controls the search strategy. The agent should prioritize trying several different configurations before refining any one of them, or should immediately refine the first promising result it finds.
\subsection{Step 2: choosing which configurations to run}
\label{sec:doe}

With five two-level axes, testing every combination means 32 configurations in total. Because a research loop is not perfectly deterministic, each configuration should be repeated several times. Running all 32 configurations with several repeats each was outside the scope of this project. Therefore, a eight-configuration Plackett--Burman design was used instead.\footnote{A specific, non-arbitrary subset of the 32 possible configurations, generated by the classical cyclic construction (a base pattern of $+,+,+,-,+,-,-$ cyclically shifted seven times, plus one row of all minus signs), chosen so that the individual effect of every one of the five axes can still be estimated without being distorted by any of the others.} The price of using only eight runs is that interactions between pairs of axes cannot be told apart from each other. This is a standard and explicitly accepted trade-off, and Section~\ref{sec:results} shows a concrete case where it matters. Table~\ref{tab:design} lists the resulting eight configurations.

\begin{table}[h]
\centering
\caption{The eight configurations used in the ablation study. Columns are the five axes from Section~\ref{sec:configspace} in order: metric, breadth, stopping rule, output format, emphasis.}
\label{tab:design}
\begin{tabular}{cccccc}
\toprule
Configuration & Metric & Breadth & Stopping & Output & Emphasis \\
\midrule
0 & 1 & 1 & 1 & 0 & 1 \\
1 & 0 & 1 & 1 & 1 & 0 \\
2 & 0 & 0 & 1 & 1 & 1 \\
3 & 1 & 0 & 0 & 1 & 1 \\
4 & 0 & 1 & 0 & 0 & 1 \\
5 & 1 & 0 & 1 & 0 & 0 \\
6 & 1 & 1 & 0 & 1 & 0 \\
7 & 0 & 0 & 0 & 0 & 0 \\
\bottomrule
\end{tabular}
\end{table}

Each configuration was run five times, with five different random seeds, giving forty runs in total, at the upper end of the three-to-five repeats that more reliably separate a genuine configuration effect from noise. Each run is an independent conversation with the language-model agent (Section~\ref{sec:limitations}), seeded only with that configuration's \texttt{program.md} and run at nonzero temperature, so the five repeats genuinely re-decide the search each time rather than replaying one plan. Both the executor's stochastic fits and the agent's own run-to-run variation contribute to the within-configuration spread that Section~\ref{sec:surrogate} relies on.
\subsubsection{What \texttt{program.md} contains}

\texttt{program.md} is comprised of a fixed statement of the research question, a fixed description of the dataset and the tool available for running an experiment, and a fixed specification of the two-shape JSON format every response must follow. Only evaluation, search strategy, budget and stopping, reporting style actually changes. Listing~\ref{lst:programmd} shows these four paragraphs exactly as generated for two configurations. Configuration 7 with vague metric, narrow, fixed budget, terse, explore-first, $(M,B,S,O,E)=(0,0,0,0,0)$ and configuration 1 with vague metric, broad, adaptive, verbose, explore-first, $(0,1,1,1,0)$. Everything outside these four paragraphs is character-for-character identical between the two files.

\begin{lstlisting}[caption={The response format specification. Identical in every one of the 32 possible \texttt{program.md} files.},label={lst:jsonformat}]
At every step, respond with a single JSON object and nothing else,
matching one of these two shapes:

Proposing an experiment:
{"action": "propose", "hypothesis": "<why you think this configuration
 will help>", "model": "random_forest" | "gradient_boosting",
 "params": {"<hyperparameter>": <value>, ...}}

Interpreting the last result and deciding what happens next:
{"action": "interpret", "interpretation": "<what the result tells you>",
 "decision": "continue" | "stop", "final_recommendation": "<only
 required when decision == stop: your final recommended model and
 hyperparameters, and why>"}
\end{lstlisting}

\begin{lstlisting}[caption={The four instructional paragraphs that actually vary.},label={lst:programmd}]
--- Configuration 7: (M,B,S,O,E) = (0,0,0,0,0) ---

## Evaluation
Evaluate model quality in whatever way you judge appropriate.

## Search strategy
Stay within a single model family per experiment. Pick either
RandomForestClassifier or GradientBoostingClassifier for the whole
session unless you have exhausted reasonable hyperparameters for it
and explicitly decide to switch.

Prioritize breadth of search: try several meaningfully different
configurations before you spend more than one experiment refining
any single one of them.

## Budget and stopping
You have a fixed budget of exactly 3 experiments in total. Run
exactly 3, then stop and report your final recommendation,
regardless of the outcome of experiment 3.

## Reporting style
Keep every response short: state the action you are taking and the
numbers involved, with minimal prose.


--- Configuration 1: (M,B,S,O,E) = (0,1,1,1,0) ---

## Evaluation
Evaluate model quality in whatever way you judge appropriate.

## Search strategy
You are free to compare RandomForestClassifier and
GradientBoostingClassifier side by side within the same session, and
to vary multiple hyperparameters of each at once.

Prioritize breadth of search: try several meaningfully different
configurations before you spend more than one experiment refining
any single one of them.

## Budget and stopping
You have a budget of up to 8 experiments. After each experiment,
decide whether to continue: stop as soon as an experiment fails to
improve the evaluation metric (defined above) by more than 0.002
over the best result so far, or when you reach 8 experiments,
whichever comes first.

## Reporting style
Explain your reasoning at every step in full sentences: what you
hypothesize, why, what you expect to observe, and how you interpret
the actual result relative to that expectation.
\end{lstlisting}

\subsubsection{What a run actually looks like}

To make the loop mechanics concrete, Listing~\ref{lst:transcript} reproduces the first two iterations of one run. Every field under \texttt{execute} is the executor's output for the stated model and hyperparameters. The agent proposes Random Forest with 100 trees and depth 5 at first, and the executor reports back a mean cross-validated accuracy of 0.962. Because configuration 1's search-strategy paragraph favors breadth before refinement, the second proposal switches model family entirely to Gradient Boosting rather than tuning the Random Forest further. Since configuration 1's stopping rule requires an improvement of more than 0.002 to continue, this single exploratory step is itself what ends the run after only two experiments. This is a result of an interaction between the breadth and stopping-rule paragraphs.

\begin{lstlisting}[caption={First two iterations of a real run under configuration 1.},label={lst:transcript}]
step 1  propose    model=random_forest, params={n_estimators: 100,
                   max_depth: 5}
        "I'll start with a Random Forest of 100 trees and moderate
         depth 5, a sensible default for a small tabular dataset,
         to establish a baseline before exploring further."
step 1  execute    cv_accuracy_mean=0.9623  cv_accuracy_std=0.0113
                   val_accuracy=0.9474
step 1  interpret  "The baseline Random Forest already reaches a
         cross-validated accuracy of about 0.962 ... this becomes
         the reference to beat."
step 1  decide     continue

step 2  propose    model=gradient_boosting, params={n_estimators: 100,
                   max_depth: 3, learning_rate: 0.1}
        "Since I am encouraged to explore breadth before refining,
         I will now try a completely different model family ..."
step 2  execute    cv_accuracy_mean=0.9448  cv_accuracy_std=0.0150
                   val_accuracy=0.9357
step 2  interpret  "Gradient Boosting underperforms the Random
         Forest baseline here (0.945 vs 0.962) ... well past the
         point of counting as an improvement."
step 2  decide     stop
        final_recommendation: "random_forest, n_estimators=100,
         max_depth=5 (cv_mean~0.962). ... the exploratory experiment
         itself triggered the stopping rule."
\end{lstlisting}

\subsection{Step 3: metrics for a run}
\label{sec:encoding}

Because each run already produces a structured record, reducing it to a set of metric is straightforward. Five metrics were chosen, and each answers a specific question about the run. They metrics fall into three groups.

The first group asks how much research was conducted, and how widely it ranged, independent of whether any of it paid off. \textbf{Number of experiments} is best possible answer to "how much research happened". Every experiment corresponds to one real model fit. \textbf{Number of distinct model families tried} matters since the breadth axis (Section~\ref{sec:configspace}) only \emph{permits} switching families. Thus, this metric is the only way to check whether that permission was actually exercised. \textbf{Search width} answers a similar question. A run that never switches models can still search a narrow or a wide slice of that model's hyperparameters.

The second group asks whether of the aforementioned led to a better fit. \textbf{Best cross-validated accuracy reached} is  one metric that measures the outcome of the research. It matters since process and outcome can come apart. For instance, an instruction file could change how an agent searches without changing how well it eventually does. It is also the metric utilized for the causal-validation check in Step 5 (Section~\ref{sec:validation}).

Instead of being about the agent at all, the third group measures whether the pipeline can be trusted. \textbf{Average length of the agent's stated hypotheses, in characters,} is a  crude proxy for verbosity. It is included for a specific reason: the output-format axis (Section~\ref{sec:configspace}) directly  instructs the agent to be either terse or fully reasoned. Therefore this is the one metric where the correct answer is known in beforehand. In case the fitted model fails in Step 4 to identify the output-format axis as the driver of this particular metric, that is a sign that something in the pipeline is broken instead of assuming that the agent is behaving unexpectedly. It functions as a calibration instrument for every other result in this report.

The result of this step is a single table with one row per run.

The same encoding script computes two further quantities that are deliberately \emph{not} modeled in Step~4, listed here so the omission is explicit rather than silent. The number of distinct hyperparameter settings tried is, in this study, identical to the number of experiments in every one of the forty runs and the improvement in cross-validated accuracy over the first experiment is near-degenerate, taking values between 0 and about 0.005 that merely re-express differences already captured by the best-accuracy metric. Both remain in \texttt{encoding/behavioral\_table.json} for inspection but are excluded from the surrogate as redundant. Full per-axis coefficient tables and the decision-tree cross-checks for all five modeled metrics.

\subsection{Step 4: a simple model of configuration and behavior}
\label{sec:surrogate}

For each of the five behavioral quantities $y$, a linear model of the form
\[
y \;=\; \beta_0 + \beta_M x_M + \beta_B x_B + \beta_S x_S + \beta_O x_O + \beta_E x_E + \varepsilon
\]
was fit by ordinary least squares, where each $x_i \in \{-1,+1\}$ codes one of the five axes (level 0 as $-1$, level 1 as $+1$) and $\varepsilon$ is the usual error term. The model is fit on all $n=40$ individual runs rather than on the eight configuration means. Because the design is balanced (every configuration has the same five repeats) and orthogonal, this changes none of the fitted coefficients: ordinary least squares on balanced replicates returns exactly the same $\beta_i$, and hence the same effects and variance shares, as least squares on the configuration means. What it changes is the residual, and therefore every standard error, $t$-statistic, $p$-value, and $R^2$: these now reflect the genuine run-to-run scatter within each configuration and are computed with $n-6=34$ residual degrees of freedom instead of two.

Coding the axes as $\pm 1$, together with the fact that the Plackett--Burman design is orthogonal, has two useful consequences. First, each $\beta_i$ can be estimated independently of the others: the fitted value of $\beta_M$, say, does not depend on which other axes happened to be varied alongside it. Second, the share of the model's total explained variance attributable to axis $i$ is simply $\beta_i^2 / \sum_j \beta_j^2$, giving a direct, additive breakdown of "how much of the variation in $y$ this axis accounts for" without any further computation.

For every coefficient, \texttt{statsmodels} additionally reports a standard error, a $t$-statistic, and a two-sided $p$-value for the null hypothesis $\beta_i = 0$ against a $t$-distribution. Because each of the forty runs is an independent model conversation (Section~\ref{sec:limitations}) every one of the five behavioral metrics genuinely varies from seed to seed, so all five repeats are true replicates and the $t$-test with $34$ residual degrees of freedom is a real, well-powered test throughout. This is worth stating because it was not always so: an earlier iteration of this study, in which the agent's proposals were authored once per configuration rather than generated live, produced near-zero within-configuration scatter for the more deterministic metrics and, with it, meaningless $p$-values. Running a genuine per-step model removes that artifact. The price is that the run-to-run scatter is now real, and, as Section~\ref{sec:results} shows, for one metric it turns out to swamp the configuration effect entirely.

\subsection{Step 5: checking the model with a real intervention}
\label{sec:validation}

A model that only fits the data used to build it has not yet demonstrated anything beyond correlation. The last step is meant to turn the fitted model into a claim that can actually be tested. One configuration, $(0,0,0,0,1)$, was chosen because it was not among the eight used to fit the model in Step 4. Using the fitted model, two of its behavioral metrics were predicted before this configuration was ever run. The corresponding instruction file was then generated and executed five times, using five random seeds not used anywhere else in this study, and the predictions were compared against what actually happened. This is the difference between saying that the model fits the data it was trained on and the model's implied claim held up when actually tried.

\section{Experimental Setup}
\label{sec:setup}

Every run used the Breast Cancer Wisconsin diagnostic dataset as bundled with scikit-learn: 569 samples, 30 numeric features, a binary classification target. A single 70/30 train/validation split, fixed with the same random seed across every run, ensured that only the instructions and the loop's own stochastic behavior varied between configurations. The two model families available to the agent were scikit-learn's \texttt{RandomForestClassifier} and \texttt{GradientBoostingClassifier}. For the former the agent could vary the number of trees, tree depth, minimum leaf size, and the fraction of features considered at each split, and for the latter the number of boosting stages, tree depth, learning rate, and subsample fraction. Regardless of what a given configuration's instructions emphasized, the executor always reported the same measurements back: five-fold cross-validated accuracy with its standard deviation on the training split, and accuracy on the held-out validation split.

The research agent on every run was OpenAI's \texttt{gpt-4o-mini}, called fresh once per run through the loop harness (\texttt{loop/agent\_runner\_llm.py}): it saw only that run's \texttt{program.md}, proposed one experiment at a time as JSON, read back the executor's real numbers, and decided when to stop. In total, the study consisted of eight ablation configurations run five times each, plus one held-out configuration also run five times, for forty-five runs and 181 individual model-fitting experiments across all of them combined. All code, generated instruction files, and raw run records are kept alongside this report so the pipeline can be re-run or inspected directly.

\section{Results}
\label{sec:results}

Table~\ref{tab:surrogate} summarizes, for each of the five behavioral metrics from Step 3, which axis the fitted linear model identifies as the dominant driver, how large a share of the variation it accounts for, the two-sided $p$-value for that axis's coefficient, and the model's overall fit ($R^2$) across the forty runs. Full coefficient tables, with standard errors, for the two quantities discussed in most depth below follow in Tables~\ref{tab:coef-accuracy} and \ref{tab:coef-verbosity}.

\begin{table}[h]
\centering
\small
\caption{Which axis explains which behavioral quantity, according to the fitted linear model ($n=40$ runs of the language-model agent, df$_{\text{resid}}=34$). $p$ is for the dominant axis's coefficient. The four \emph{process} metrics are each cleanly driven by the axis one would expect. The one \emph{outcome} metric, best accuracy, is not driven by any axis.}
\label{tab:surrogate}
\begin{tabular}{p{3.3cm}p{1.6cm}p{1.3cm}p{1.1cm}p{4.2cm}p{0.7cm}}
\toprule
Behavioral metrics & Dominant axis & Var.\ share & $p$ & Direction of effect & $R^2$ \\
\midrule
Verbosity (avg.\ proposal length) & Output format & 94\% & $6\times10^{-12}$ & full explanations $\to$ roughly 1.8$\times$ longer proposals & 0.77 \\
Number of experiments & Stopping & 95\% & $2\times10^{-8}$ & adaptive budget $\to$ more experiments ($+2.4$) & 0.62 \\
Model families tried & Breadth & 61\% & $2\times10^{-6}$ & broad permission $\to$ more families ($+0.6$) & 0.61 \\
Search width & Emphasis & 63\% & 0.005 & exploit-first $\to$ narrower spread ($-55$) & 0.30 \\
Best accuracy reached & \emph{none} & --- & $\ge 0.45$ & no axis moves it. Agent lands near 0.96 regardless & 0.03 \\
\bottomrule
\end{tabular}
\end{table}

The first row is, by design, close to a sanity check. The output-format axis directly instructs the agent to write either terse or fully reasoned responses, so finding that it explains 94\% of the variation in how long the agent's stated hypotheses are is exactly what should happen if the whole pipeline is working correctly end to end. Table~\ref{tab:coef-verbosity} gives the full picture: all five coefficients, their standard errors, and their $p$-values. The output-format coefficient dwarfs every other -- a level-0-to-1 change lengthens the average proposal by about 125 characters (from roughly 146 to 266), against effects of at most about twenty characters for the other four axes, none of which is significant. Crucially, the fit is good but not perfect ($R^2=0.77$): the agent's proposal lengths do vary from run to run even at a fixed configuration, because these are genuinely independent model calls, and the output axis nonetheless captures the overwhelming majority of the variation, at $p=6\times10^{-12}$. This is a real, well-powered result, not the pseudoreplicated artifact the earlier authored-proposal version produced. If some unrelated axis had turned out to carry that variance share instead, it would have pointed to a bug in the pipeline rather than a finding. Because this check passes cleanly on a genuine agent, it is reasonable to trust the pipeline on the less obvious results that follow including the negative one.

\begin{table}[h]
\centering
\caption{Full linear-model coefficients for verbosity (average proposal length, in characters). Coefficient is on the $\pm1$-coded axis. A level-0-to-1 change is twice this value. $n=40$, df$_{\text{resid}}=34$, $R^2=0.77$. Output format is the only significant axis and carries 94\% of the variance share. The standard error is identical across axes ($\sqrt{\text{MSE}/40}$) because the $\pm1$ design is balanced and orthogonal ($X^\top X = 40\,I$). It differs between tables only through each metric's residual variance.}
\label{tab:coef-verbosity}
\begin{tabular}{lrrrr}
\toprule
Axis & Coefficient & Std.\ error & $t$ & $p$ \\
\midrule
Metric (M) & $+10.03$ & 6.07 & $1.65$ & 0.108 \\
Breadth (B) & $+3.04$ & 6.07 & $0.50$ & 0.620 \\
Stopping (S) & $-11.29$ & 6.07 & $-1.86$ & 0.072 \\
Output format (O) & $+62.29$ & 6.07 & $10.3$ & $\mathbf{6\times10^{-12}}$ \\
Emphasis (E) & $-2.41$ & 6.07 & $-0.40$ & 0.694 \\
\bottomrule
\end{tabular}
\end{table}

The most important result is a negative one. The best cross-validated accuracy a run reaches is essentially not controlled by any of the five instructional axes. The linear model explains almost none of the run-to-run variation in it ($R^2=0.03$), no axis's coefficient comes anywhere near significance (every $p\ge0.45$. Table~\ref{tab:coef-accuracy}), and the largest effect any axis has is about a fifth of an accuracy percentage point, smaller than the ordinary seed-to-seed scatter within a single configuration. Whatever the instructions say about metric, breadth, stopping, output, or search emphasis, the agent converges on a cross-validated accuracy of roughly 0.96 on this dataset. 

This directly overturns the headline finding of an earlier version of this study. In that version the agent's proposals were authored in advance, once per configuration, rather than generated live. On that data the emphasis axis came out as the dominant driver of accuracy, credited with 80\% of the between-configuration variance. That effect did not survive replacing the authored plan with a genuine per-step model: it was a property of how the proposals had been written, not of an autoresearch loop. The honest reading is that, for this task, the instruction file shapes \emph{how the agent works} but not \emph{how well it does}. The outcome is pinned near the dataset's ceiling regardless. Catching exactly this kind of false positive is what moving from a scripted stand-in to a real agent was for.

\begin{table}[h]
\centering
\caption{Full linear-model coefficients for best cross-validated accuracy reached. Coefficient is on the $\pm1$-coded axis. A level-0-to-1 change is twice this value. $n=40$, df$_{\text{resid}}=34$, $R^2=0.03$. No axis is significant (every $p\ge0.45$): configuration explains essentially none of the run-to-run variation in accuracy. The standard error is equal across axes by the design's orthogonality.}
\label{tab:coef-accuracy}
\begin{tabular}{lrrrr}
\toprule
Axis & Coefficient & Std.\ error & $t$ & $p$ \\
\midrule
Metric (M) & $+0.000252$ & 0.00049 & $0.52$ & 0.609 \\
Breadth (B) & $-0.000127$ & 0.00049 & $-0.26$ & 0.796 \\
Stopping (S) & $+0.000374$ & 0.00049 & $0.77$ & 0.450 \\
Output format (O) & $+0.0000008$ & 0.00049 & $0.00$ & 0.999 \\
Emphasis (E) & $-0.0000022$ & 0.00049 & $-0.00$ & 0.996 \\
\bottomrule
\end{tabular}
\end{table}

The four \emph{process} rows, by contrast, come out clean, well-powered, and each driven by the axis one would expect. The number of experiments run is dominated by the stopping-rule axis ($95\%$ of the variance, $p=2\times10^{-8}$): fixed-budget runs execute their three experiments, while under the adaptive rule the model keeps going, averaging closer to five or six before it judges the returns exhausted, so switching that axis on adds about $2.4$ experiments. How many distinct model families get tried is led, unsurprisingly, by the breadth axis ($61\%$, $p=2\times10^{-6}$). The axis that permits or forbids trying both Random Forest and Gradient Boosting at all, with the metric and emphasis axes as minor secondary influences. Search width, the range of tree counts spanned, is led by the emphasis axis ($63\%$, $p=0.005$): an agent told to exploit the first promising result immediately stays in a narrower neighborhood of hyperparameters than one told to explore first. In every one of these three cases the dominant axis is exactly the paragraph of \texttt{program.md} that names the behavior in question, and each is now a genuine, well-powered result computed over independent runs.

The held-out configuration used for causal validation was deliberately excluded from the eight configurations used to fit the model above. Table~\ref{tab:validation} compares what the fitted model predicted for two of its behavioral quantities, before this configuration was ever executed, against what was actually observed once it was run five times with five fresh random seeds.

\begin{table}[h]
\centering
\caption{Predicted versus observed behavior for the held-out configuration.}
\label{tab:validation}
\begin{tabular}{lccc}
\toprule
Behavioral quantity & Predicted & Observed & Difference \\
\midrule
Number of experiments & 2.95 & 3.00 & $+0.05$ \\
Best accuracy reached & 0.9609 & 0.9633 & $+0.0024$ \\
\bottomrule
\end{tabular}
\end{table}

The two predictions land very differently, and that difference is itself the point. The number-of-experiments prediction is essentially exact: the model had correctly identified the stopping-rule axis as the driver of how much research happens, and for this fixed-budget held-out configuration it predicted $2.95$ against an observed $3.00$, which is a miss of one-twentieth of an experiment. This is a real, checkable claim about a run that had not happened when the prediction was made, and it held. The accuracy prediction, by contrast, is uninformative, and honestly so: the accuracy surrogate explains essentially none of the variation in the first place ($R^2=0.03$), so its point prediction of $0.9609$ against an observed $0.9633$ is not evidence of anything. Both numbers are just about 0.96, which is where every configuration lands. The validation step thus does two useful jobs at once. It confirms the one part of the model that carries real signal, and it refuses to dress up the part that does not: a surrogate with $R^2=0.03$ makes no prediction worth trusting, and the held-out check says so rather than papering over it.

\section{How much of the black box has actually been opened}
\label{sec:blackbox}

The introduction described the problem this project set out to address as a black box: change \texttt{program.md} and the agent's behavior changes, but there was no way to say which part of the instructions was responsible for which part of the change, nor to predict what a change not yet tried would do. It is worth stating precisely what opening that black box meant in practice, what evidence from Sections~\ref{sec:methodology} and~\ref{sec:results} actually supports that claim, and what kind of black box this project did not open, so that the white-box claim in the title of this line of work is not read as broader than it is.

\paragraph{What became visible.} Three specific things changed between "we observe that the instructions matter" and where this study ends up. First, the input stopped being an undifferentiated block of text and became a small, labeled vector: five named axes, each with a stated meaning, rather than one opaque file. This alone is a real step and it is what made everything downstream possible. Second, the relationship between that vector and the agent's behavior was captured in a model whose insides can be read directly: a linear equation with five coefficients, not a black-box predictor of its own. Every claim in Table~\ref{tab:surrogate} can be traced back to one number in one regression, and Tables~\ref{tab:coef-verbosity} and~\ref{tab:coef-accuracy} show all five of those numbers, not just the largest one. Third, and most importantly, one specific claim generated by that model was checked against a real intervention before the outcome was known (Section~\ref{sec:validation}), rather than only checked against the data used to build the model. A model that merely correlates with data it was fit on can always be produced. A model whose stated prediction is checked against a run that had not happened yet has done something a black box by definition cannot do: made a falsifiable claim and been tested on it.

\paragraph{What stayed exactly as opaque as before.} None of this says anything about \emph{why} a large language model responds to instructions the way it does at the level of what actually happens inside it. There is no claim here about attention patterns, learned representations, or any mechanism internal to the model. The language model itself remains as much of a black box after this report as before it. What changed is that its \emph{external, input-output behavior}, for one narrow task and one small set of instructional choices, now has a description that is interpretable and testable, in the same sense that a linear regression relating a drug's dosage to a patient outcome is interpretable without anyone needing to understand human biochemistry to use it. This is precisely the sense in which AutoML hyperparameter-importance tools "open" the black box of a learning algorithm (Section~\ref{sec:background}): functional ANOVA over a Random Forest's hyperparameters tells you that tree depth matters more than leaf size for a given dataset without telling you anything about how decision trees or bagging actually work internally. The white box here is a white box at the level of \emph{behavior}, not at the level of \emph{mechanism}, and the two should not be confused.

\paragraph{Where the white-box claim is, and is not, entitled to travel.} Even restricted to behavior, the claim earned by this study is narrower than it might sound. It is entitled to say: for this task, this dataset, these five axes, and this one language model, the relationship between instruction and behavior can be measured, summarized in an interpretable model, and used to make a checkable prediction that at least partially held up. It is not entitled to say that this relationship holds for other tasks, other datasets, a larger or differently chosen set of axes, or for language models in general rather than for the one small model (\texttt{gpt-4o-mini}) that acted as the agent here. Every one of the specific numbers in Section~\ref{sec:results} (the 94\%, the near-zero accuracy $R^2$, the $p$-values) describes this study's own small, self-contained system, not autoresearch loops in general. What generalizes, if anything does, is the method: the demonstration that an instruction file can be treated as a configuration space, ablated rather than randomly sampled, modeled with something interpretable, and checked with a real intervention, at all. Whether the specific findings generalize is exactly the open, unanswered question the next steps in Section~\ref{sec:conclusion} are aimed at.

To put the whole claim in one sentence: this project turned a black box into a white box in the narrow, specific sense of turning "the instructions matter, somehow" into "here is a model of which instruction changes which behavior, by how much, and it made one falsifiable prediction that was checked",and it did this for one small system standing in for the real thing, not for autoresearch loops as a general phenomenon.

\section{Limitations}
\label{sec:limitations}

The way the agent's decisions were produced is now largely resolved, and the resolution is what makes the results above worth taking seriously. Each of the forty-five runs was a genuinely independent call to a language model (OpenAI \texttt{gpt-4o-mini}) that saw nothing but that run's \texttt{program.md}, with no memory of any other run, deciding each next experiment live from the real numbers it read back. That is the design the study always pointed at, and reaching it is precisely what exposed the accuracy finding of the earlier, authored-proposal version as an artifact.

What remains is a limitation of a different kind: the agent is one specific, small model. Everything in Section~\ref{sec:results} is a property of \texttt{gpt-4o-mini} on this task, not of autoresearch loops in general. A larger or differently trained model might search more aggressively, might obey the instruction paragraphs more or less faithfully, and might well show accuracy effects this model does not. Re-running the identical design across several models, and checking which of these effects are stable across them, is now the single most useful next step. A related caveat is that the accuracy null may be partly a property of the \emph{task}: the breast-cancer dataset is easy enough that almost any reasonable tree-ensemble configuration reaches about 0.96, so there is little headroom for an instructional choice to move the outcome even in principle. A harder task might restore one.

A second limitation is scale: eight configurations support only individual-axis effects, not interactions between pairs of axes. Estimating those interactions directly would require going beyond the eight Plackett--Burman configurations, for example to a resolution-IV or full factorial design. A third and final limitation is scope: every experiment concerns one small dataset and two model families with a modest set of hyperparameters, so the specific numbers reported here should not be assumed to generalize beyond this setting without repeating the study elsewhere.

\section{Conclusion}
\label{sec:conclusion}

This study carried out all five planned steps toward explaining, rather than merely observing, how an instruction file shapes the behavior of an automated research loop, by treating the instruction file the same way automated machine learning treats a hyperparameter vector: as a small set of discrete choices whose individual effect on behavior can be measured with a fractional factorial design, summarized with a simple interpretable model, and checked with a real held-out intervention. Every number in this report came from actually executed machine learning experiments, with an independent language model (\texttt{gpt-4o-mini}) driving each run live from nothing but its instruction file, so the effects reported are properties of a real agent, not of a hand-written stand-in.

The method recovered the expected relationships cleanly: the four axes that describe how the agent should \emph{work} each turned out to be the dominant driver of exactly the behavior it names, with well-powered statistics over independent runs. Its most informative result, though, was a negative one that only appeared once a real model replaced the scripted stand-in: none of the five instructional choices reliably moves the accuracy the agent reaches, which sits near the dataset's ceiling however the instructions are worded. An earlier version of this study, with authored proposals, had confidently reported the opposite. The held-out validation step reflected the correction honestly, confirming the process prediction it had signal for and declining to trust the outcome prediction it did not.

The most useful next steps, in order, are: repeating the identical design across several different language models, to see which of these effects are stable across models rather than specific to \texttt{gpt-4o-mini}. Expanding the ablation design beyond eight configurations so that interactions between axes can be estimated directly, and repeating the study on harder tasks with real headroom in the outcome, to test whether the accuracy null is a property of the agent or merely of an easy benchmark.

\appendix

\section{Full surrogate coefficients for all five metrics}
\label{sec:appendix-coef}

Section~\ref{sec:results} gives full per-axis coefficient tables only for the two metrics discussed in depth (verbosity and best accuracy, Tables~\ref{tab:coef-verbosity} and~\ref{tab:coef-accuracy}) and summarizes the remaining three in Table~\ref{tab:surrogate}. Table~\ref{tab:appendix-coef} collects the complete linear-model output for all five modeled metrics, so that every number behind Table~\ref{tab:surrogate} is available in one place. All models are fit on the $n=40$ individual runs (df$_{\text{resid}}=34$). Coefficients are on the $\pm1$-coded axes, so a level-0-to-1 change equals twice the coefficient. As explained in Section~\ref{sec:surrogate}, the standard error is identical across the five axes within each metric because the design is balanced and orthogonal ($X^\top X = 40\,I$), and differs between metrics only through each metric's residual variance.

\begin{table}[htbp]
\centering
\small
\caption{Complete linear main-effects coefficients for all five behavioral metrics ($n=40$ runs of the language-model agent, df$_{\text{resid}}=34$). ``Var.\ share'' is $\beta_i^2/\sum_j\beta_j^2$. Every run is an independent model call, so within-configuration scatter is genuine and all $p$-values are honest. For best accuracy, note $R^2=0.03$ and that no axis is significant.}
\label{tab:appendix-coef}
\begin{tabular}{lrrrr r}
\toprule
Axis & Coefficient & Std.\ error & $t$ & $p$ & Var.\ share \\
\midrule
\multicolumn{6}{l}{\textit{Best CV accuracy reached} \quad ($R^2=0.03$)}\\
\quad Metric (M)        & $+0.000252$  & 0.000489 & $0.52$  & 0.609 & 28.9\% \\
\quad Breadth (B)       & $-0.000127$  & 0.000489 & $-0.26$ & 0.796 & 7.4\% \\
\quad Stopping (S)      & $+0.000374$  & 0.000489 & $0.77$  & 0.450 & 63.7\% \\
\quad Output (O)        & $+0.0000008$ & 0.000489 & $0.00$  & 0.999 & 0.0\% \\
\quad Emphasis (E)      & $-0.0000022$ & 0.000489 & $-0.00$ & 0.996 & 0.0\% \\
\midrule
\multicolumn{6}{l}{\textit{Number of experiments} \quad ($R^2=0.62$)}\\
\quad Metric (M)        & $-0.15$ & 0.164 & $-0.92$ & 0.366 & 1.5\% \\
\quad Breadth (B)       & $-0.10$ & 0.164 & $-0.61$ & 0.546 & 0.7\% \\
\quad Stopping (S)      & $+1.20$ & 0.164 & $7.32$  & $\mathbf{2\times10^{-8}}$ & 95.0\% \\
\quad Output (O)        & $+0.20$ & 0.164 & $1.22$  & 0.231 & 2.6\% \\
\quad Emphasis (E)      & $-0.05$ & 0.164 & $-0.31$ & 0.762 & 0.2\% \\
\midrule
\multicolumn{6}{l}{\textit{Distinct model families tried} \quad ($R^2=0.61$)}\\
\quad Metric (M)        & $-0.15$ & 0.052 & $-2.88$ & \textbf{0.007} & 15.3\% \\
\quad Breadth (B)       & $+0.30$ & 0.052 & $5.75$  & $\mathbf{2\times10^{-6}}$ & 61.0\% \\
\quad Stopping (S)      & $-0.10$ & 0.052 & $-1.92$ & 0.064 & 6.8\% \\
\quad Output (O)        & $-0.05$ & 0.052 & $-0.96$ & 0.345 & 1.7\% \\
\quad Emphasis (E)      & $-0.15$ & 0.052 & $-2.88$ & \textbf{0.007} & 15.3\% \\
\midrule
\multicolumn{6}{l}{\textit{Search width (n\_estimators spread)} \quad ($R^2=0.30$)}\\
\quad Metric (M)        & $-2.5$  & 9.13 & $-0.27$ & 0.786 & 0.5\% \\
\quad Breadth (B)       & $+5.0$  & 9.13 & $0.55$  & 0.587 & 2.1\% \\
\quad Stopping (S)      & $+20.0$ & 9.13 & $2.19$  & \textbf{0.035} & 33.5\% \\
\quad Output (O)        & $+2.5$  & 9.13 & $0.27$  & 0.786 & 0.5\% \\
\quad Emphasis (E)      & $-27.5$ & 9.13 & $-3.01$ & \textbf{0.005} & 63.4\% \\
\midrule
\multicolumn{6}{l}{\textit{Verbosity (avg.\ proposal chars)} \quad ($R^2=0.77$)}\\
\quad Metric (M)        & $+10.03$ & 6.07 & $1.65$  & 0.108 & 2.4\% \\
\quad Breadth (B)       & $+3.04$  & 6.07 & $0.50$  & 0.620 & 0.2\% \\
\quad Stopping (S)      & $-11.29$ & 6.07 & $-1.86$ & 0.072 & 3.1\% \\
\quad Output (O)        & $+62.29$ & 6.07 & $10.3$  & $\mathbf{6\times10^{-12}}$ & 94.1\% \\
\quad Emphasis (E)      & $-2.41$  & 6.07 & $-0.40$ & 0.694 & 0.1\% \\
\bottomrule
\end{tabular}
\end{table}



\begin{thebibliography}{9}

\bibitem{fawcetthoos}
C. Fawcett and H. H. Hoos.
Analysing differences between algorithm configurations through ablation.
Journal of Heuristics, 22(4):431--458, 2016.

\bibitem{hutter2014fanova}
F. Hutter, H. Hoos, and K. Leyton-Brown.
An efficient approach for assessing hyperparameter importance.
Proceedings of the 31st International Conference on Machine Learning (ICML), 2014.

\bibitem{plackettburman}
R. L. Plackett and J. P. Burman.
The design of optimum multifactorial experiments.
Biometrika, 33(4):305--325, 1946.

\end{thebibliography}

\end{document}

